\chapter{Minería: encontrar patrones}
\label{cap:s-patrones}
\textit{Materia: Minería de datos. Pregunta: ¿qué patrones de tiempo, de lugar y de comportamiento explican dónde se
concentra la gente?}

La minería de datos busca lo que no se ve a simple vista. En este proyecto se buscaron seis patrones; cada uno alimenta
una decisión de la campaña.

\section{¿Qué tan concentrado está el turismo?}
\begin{intuicion}
Imagina un pastel repartido entre varias personas. Si una se lleva casi todo, el reparto está concentrado. Hay un número,
el \textbf{índice de Herfindahl-Hirschman}, que resume qué tan desigual es el reparto: vale 0 si todos reciben lo mismo y
1 si uno se lleva todo. Es el que usan las autoridades de competencia para ver si un mercado está acaparado.
\end{intuicion}

\begin{ejemplo}[ · los pasajeros de avión en 2024]
De 15,959,277 pasajeros que llegaron en avión al estado, Cancún recibió 14,769,302: el \textbf{92.5\,\%}. Tulum, 3.9\,\%;
Cozumel, 2.2\,\%; Chetumal, 1.4\,\%.

El índice eleva al cuadrado la parte de cada aeropuerto y las suma: $0.9254^2+0.0389^2+0.0221^2+0.0136^2=0.8586$. Con
cuatro aeropuertos, lo más parejo posible sería que cada uno tuviera la cuarta parte, y el índice valdría 0.25. Se ajusta
para que ese caso valga 0:
\[
\frac{0.8586-0.25}{1-0.25}=\frac{0.6086}{0.75}=\mathbf{0.811}.
\]
Un 0.811 en una escala de 0 a 1: casi todo entra por una sola puerta.
\end{ejemplo}

\begin{center}
\small
\begin{tabular}{lrr}
\encabezado \h{Qué se mide} & \h{Concentración (0 a 1)} & \h{Parte de los 5 lugares} \\
Pasajeros de avión (2024) & 0.811 & 1.4\,\% \\
Cuartos de hotel (julio de 2026) & 0.219 & 1.7\,\% \\
Visitantes a zonas arqueológicas (2025) & 0.276 & 4.1\,\% \\
Negocios turísticos (DENUE) & 0.133 & 14.4\,\% \\
Población (Censo 2020) & 0.225 & 12.3\,\% \\
\bottomrule
\end{tabular}
\normalsize
\end{center}
\textbf{Lectura}: los negocios están donde vive la gente (14.4\,\% contra 12.3\,\%), pero los visitantes no (entre 1.4\,\% y
4.1\,\%). \textbf{Para la campaña}: el sur tiene capacidad ociosa; promoverlo no crea una saturación nueva, aprovecha lo
que ya existe.

\section{¿Qué meses se llenan? La forma del año}
\begin{intuicion}
Para saber si enero es un mes ``alto'', se compara enero con el promedio de su propio año. Si un año tuvo 5,000
visitantes al mes en promedio y enero tuvo 8,000, enero fue 1.6 veces un mes normal. Se hace lo mismo con varios años y
se promedia. El resultado es la \textbf{forma del año}: un número por mes, donde 1 es un mes normal.
\end{intuicion}

\begin{ejemplo}[ · enero en la Ruta de las pirámides]
En 2019, Kohunlich y Dzibanché recibieron 64,139 visitantes: $64{,}139\div12=5{,}344.9$ al mes en promedio. En enero
llegaron 7,935: $7{,}935\div5{,}344.9=1.48$. Enero de 2019 fue 1.48 veces un mes normal.

Con los seis años completos que hay (2016 a 2019, 2022 y 2023), las razones de enero fueron 1.55, 1.71, 1.83, 1.48, 1.25
y 1.82. Su promedio es \textbf{1.606}: en enero llega \textbf{61\,\% más} gente que en un mes promedio. En septiembre, el
número es 0.496: \textbf{la mitad}.
\end{ejemplo}

\begin{center}
\small
\begin{tabular}{lll}
\encabezado \h{Lugar} & \h{Mes más lleno} & \h{Mes más vacío} \\
Ruta de las pirámides & enero (1.61) & septiembre (0.50) \\
Bahía Calderitas–Oxtankah & diciembre (1.40) & septiembre (0.65) \\
Chetumal (cruces de Belice) & diciembre (1.17) & febrero (0.87) \\
Cancún (referencia) & marzo (1.08) & septiembre (0.87) \\
\bottomrule
\end{tabular}
\normalsize
\end{center}

\textbf{Qué tanto pesa la temporada}: en la Ruta, el 79\,\% de lo que sube y baja un mes a otro se explica por la
temporada; el resto es azar. Es una temporada fuerte y predecible.

\textbf{Dos hallazgos}: (1) las zonas del sur suben y bajan al mismo tiempo que Cancún, pero Chetumal con Belice no;
(2) septiembre es el mes más vacío casi en todos lados y coincide con el pico de tormentas.

\begin{porque}
\textbf{Multiplicar en lugar de sumar.} Se podía decir ``enero tiene 2,000 visitantes más'' o ``enero tiene 61\,\% más''.
Kohunlich reabrió en 2025 con la mitad de la gente que tenía en 2019: ``2,000 más'' ya no tiene sentido, pero ``61\,\%
más'' sí. Por eso la forma del año se calcula en proporciones.

\textbf{Solo años completos.} El método más conocido (STL) necesita una serie sin huecos, y estas los tienen: la pandemia
y los cierres por obras. Con STL se habrían tirado años enteros. Se usó como segunda opinión donde sí se puede, y coincide
entre 85\,\% y 97\,\%.
\end{porque}

\textbf{Para la campaña}: no se anuncia un lugar en sus meses más llenos y se empuja en sus meses con espacio.

\section{¿Qué lugar está más presionado? El índice de presión}
Cada lugar se mide con cosas distintas: Chetumal tiene Tren Maya, la Ruta tiene visitantes a sus zonas, Cancún tiene
ocupación hotelera. Para compararlos se necesita una sola escala.

\begin{intuicion}
Es como calificar a alumnos de grupos distintos con exámenes distintos. Primero se pasa cada examen a una escala de 0 a 1
(0 = la calificación más baja de todo el estado, 1 = la más alta) y luego se promedian los exámenes que presentó cada
alumno. Así sale el \textbf{índice de presión turística}: 0 es lo más tranquilo y 1 lo más presionado que se ha visto en el
estado.
\end{intuicion}

\begin{ejemplo}[ · Cancún, julio de 2026]
Cancún tiene cuatro medidas ese mes. Cada una se pasa a la escala de 0 a 1 dividiendo lo que tuvo entre el rango de todo
el estado:
\begin{itemize}
\item Ocupación hotelera: tuvo 68.78\,\%; en el estado va de 17.97\,\% a 88.40\,\%. Escala:
  $(68.78-17.97)\div(88.40-17.97)=50.81\div70.43=0.721$.
\item Gente que bajó del Tren Maya por cada mil habitantes: 25.1; el máximo del estado es 505.2. Escala: 0.050.
\item Visitantes a sus zonas arqueológicas por cada mil habitantes: casi nada frente al máximo. Escala: 0.0001.
\item Llegadas por cuarto de hotel: 0.47 frente a un máximo de 432. Escala: 0.001.
\end{itemize}
Promedio: $(0.721+0.050+0.0001+0.001)\div4=0.772\div4=\mathbf{0.193}$.
\end{ejemplo}

Después se ordenan todos los lugares y meses del estado y se ponen dos cortes: la mitad más baja es
\textbf{tranquilo}, el 10\,\% más alto es \textbf{saturado} y lo de en medio es \textbf{concurrido}. Con los cortes del
índice que publica el Radar (0.186 y 0.756), Cancún, con 0.193, queda \textbf{concurrido} en julio de 2026.

\textbf{Julio de 2026}: los cinco lugares de la campaña están \textbf{tranquilos} (la Ruta con 0.129, Chetumal 0.025,
Oxtankah 0.020 y Maya Ka'an 0.009). Están concurridos Mahahual (0.731, por los cruceros), Isla Mujeres, Bacalar, Holbox,
Cancún y Cozumel. La Laguna Milagros no tiene ninguna medida: sale ``sin dato oficial'', nunca ``tranquila''.

\begin{porque}
\textbf{Pesos iguales.} Ninguna fuente dice que el tren pese más que los cruceros. Unos pesos inventados serían difíciles
de defender. Para comprobar que no cambia el resultado, se movió cada peso a la mitad y al 150\,\%: el estado cambia en
6\,\% de los casos como máximo.

\textbf{Cortes comunes para todo el estado.} Si cada lugar se comparara solo contra sí mismo, todos saldrían ``saturados''
una parte del tiempo, aunque nunca se llenen. Con la misma vara, el norte y el sur se pueden comparar.
\end{porque}

\begin{error}
\textbf{El índice que decía que Cancún estaba tranquilo.} La primera versión usaba solo los datos del gobierno del estado.
Como el estado dejó de publicar la ocupación en 2025, Cancún se quedó solo con medidas casi en cero y salió
``tranquilo'' en julio de 2026, con 0.025. SECTUR decía que ese mes tenía 68.8\,\% de ocupación. Un índice que dice que
Cancún está tranquilo en julio no sirve. Se compararon cuatro arreglos con datos, y Brandon eligió el que usa la ocupación
de SECTUR para los cuatro lugares del norte que mide, y deja fuera las medidas que solo tiene un lugar (como Belice, que
solo tiene Chetumal y se comparaba contra sí mismo).
\end{error}

\section{¿A qué se parece el norte en el país? Agrupamiento}
Cada centro turístico del país tiene una ``huella'' de 12 números: qué tan lleno está cada mes del año. Un método de
\textbf{agrupamiento} (el de Ward) junta los centros con huellas parecidas: empieza con cada centro solo y va uniendo los
más parecidos, de dos en dos, como un árbol genealógico al revés.

¿Cuántos grupos? Se mide qué tan bien acomodado está cada centro en su grupo (la \emph{silueta}: de $-1$, mal acomodado, a
1, perfecto) y se eligen los grupos con el mejor promedio. Con 55 centros, lo mejor fueron \textbf{2 grupos} (silueta
0.450):
\begin{itemize}
\item \textbf{Grupo 1}, playas muy llenas (64.5\,\% de ocupación en promedio): Cancún, Riviera Maya, Playacar, Akumal, Playa
  del Carmen, Los Cabos.
\item \textbf{Grupo 2}, el resto del país (41.1\,\%), incluidos Cozumel e Isla Mujeres.
\end{itemize}
\textbf{Para la campaña}: la saturación del norte no es una impresión; está entre los destinos más llenos de México.

\section{¿Qué semanas tienen clima raro? Isolation Forest}
\begin{intuicion}
Imagina puntos en una hoja y que los separas con cortes de tijera al azar. Un punto que está lejos de todos queda solo
con muy pocos cortes; uno en medio de la multitud necesita muchísimos. \textbf{Isolation Forest} hace eso con 300
``árboles'' de cortes al azar y mide cuántos cortes tarda en aislar cada semana. Las semanas raras se aíslan rápido.
\end{intuicion}

Cada semana se describe con su lluvia, la lluvia del peor día, el viento, el calor y la época del año (para que un
aguacero de septiembre no salga raro solo por ser septiembre). El método da un puntaje de 0 a 1: cerca de 0.5 es normal;
cerca de 1, muy raro.

\begin{ejemplo}[ · Chetumal, semana de la tormenta Nadine (octubre de 2024)]
Llovieron 148.7 mm en la semana, 65.4 mm el peor día. El bosque tardó unos 5 cortes en aislarla, cuando una semana normal
tarda unos 10. Su puntaje fue \textbf{0.712}; el corte para ``raro'' era 0.543. Como $0.712>0.543$, la semana es rara.
\end{ejemplo}

\textbf{La prueba de que funciona}: el bosque marcó como raras las tres semanas en que pasaron tormentas (Lisa en 2022,
Nadine y Sara en 2024) \textbf{sin saber que hubo tormenta}. Nadie se lo dijo.

\textbf{Para la campaña}: una semana rara pausa el anuncio de ese lugar (capítulo~\ref{cap:s-torre}).

\section{¿Qué molesta y qué enamora al turista? Las reseñas}
Se analizaron \textbf{85,987 opiniones} de turistas en Quintana Roo (Tulum, Isla Mujeres y Bacalar; ninguna de los cinco
lugares, así que sirven para saber qué molesta en los lugares llenos, no para opinar del sur).

\textbf{Temas con una lista de palabras a la vista.} Cada tema es una lista de palabras que se puede leer: ``ruido''
incluye \emph{ruido, escándalo, música alta, fiesta}; ``calma'' incluye \emph{tranquilo, paz, relajado, silencio}. Se
eligió así, y no con un modelo automático de temas, porque cada tema se puede defender palabra por palabra.

\begin{ejemplo}[ · el ruido]
De las 85,987 reseñas, 3,597 son malas (1 o 2 estrellas): $3{,}597\div85{,}987=4.18\,\%$.

De las 2,038 que hablan de ruido, 235 son malas: $235\div2{,}038=11.53\,\%$.

$11.53\div4.18=\mathbf{2.76}$. Una reseña que habla de ruido es mala \textbf{2.76 veces} más seguido de lo normal. A esto se le
llama \emph{riesgo relativo}, y es la misma medida que usan los médicos para comparar riesgos.
\end{ejemplo}

\begin{center}
\small
\begin{tabular}{lrl}
\encabezado \h{Tema} & \h{Riesgo relativo} & \h{Qué significa} \\
Ruido & 2.76 & Hunde la reseña \\
Suciedad & 1.89 & Hunde \\
Precio & 1.66 & Hunde \\
Multitudes & 1.52 & Hunde \\
Servicio & 1.18 & Casi neutro \\
Naturaleza & 0.73 & Protege \\
Cultura & 0.48 & Protege \\
Calma & 0.41 & Protege mucho \\
\bottomrule
\end{tabular}
\normalsize
\end{center}

\textbf{Combinaciones} (método Apriori, que busca qué temas aparecen juntos): una reseña que habla de ruido \emph{y} de mal
servicio es mala 2.82 veces más seguido de lo normal.

\textbf{Las palabras del turista}: las reseñas de 5 estrellas dicen \emph{excelente, increíble, hermoso, deliciosa, amable,
ruinas, ambiente}; las malas dicen \emph{peor, horrible, dinero, caro, sucia, grosero}.

\textbf{La conclusión que define la campaña}: lo que hunde una reseña (ruido, suciedad, multitudes) es lo que sobra en los
lugares llenos, y lo que la protege (calma y cultura) es justo lo que el sur ofrece. Esa es la promesa de la marca.

\section{¿Qué recomendar cerca? Clasificar negocios por su nombre}
Para que la página recomiende qué hacer, dónde comer y dónde dormir, se clasificaron los negocios del DENUE de las
localidades de los cinco lugares leyendo su \textbf{nombre} (y, si el nombre no dice nada, su giro oficial). Se hizo con
una lista de raíces de palabras, porque no había ejemplos etiquetados para entrenar un modelo.

Las reglas salieron de errores reales: ``MICHELADAS'' contiene ``HELAD'' y salía como heladería; ``BARBACOA'' contiene
``BAR''. Por eso una raíz cuenta solo al inicio de una palabra. Al revisar a mano dos muestras al azar de 40 negocios, la
segunda acertó \textbf{39 de 40}. De 2,042 negocios turísticos, 1,587 quedaron clasificados.
